% !TeX root = ./RBE550_assignment_template.tex
\documentclass{RBEassignment}

\usepackage{lipsum}






\begin{document}


\def\RBEtitle{Assignment\\Template}
\def\RBEcourse{RBE-001 Generic RBE Course}
\def\RBEsemester{Spring 2025}
\def\RBEauthor{Daniel Montrallo Flickinger, PhD}
\def\RBEnumber{0}
\def\RBEdue{2025-03-04 @ 12:00 UTC}



\RBEassignmentheader{}



\section{Background}

% TODO: write up how to use the template

% TODO: link to Github page for this project
The most recent version of the \href{https://github.com/dmflickinger/RBE550-assignment-template}{RBE Assignment Template} is on Github.  

\lipsum[2]

\subsection{2D Obstacle Fields}


\begin{figure}
\begin{center}
\begin{tikzpicture}
\coordinate (gridsize) at (15,15);
\pgfmathsetmacro{\numTets}{15}

\begin{scope}[local bounding box=L]
\foreach \i in {1,...,\numTets} {
    \pgfmathsetmacro{\tetRow}{int(12 * random())}
    \pgfmathsetmacro{\tetCol}{int(12 * random())}
    \pgfmathsetmacro{\whichTet}{int(5 * random())}
    \placeTet{\whichTet}{\tetRow}{\tetCol};
}
\draw (0,0) grid (gridsize);
\end{scope}
\end{tikzpicture}
\end{center}
\caption{Example obstacle field.%
\label{fig:obstacleField}}
\end{figure}


Figure~\ref{fig:obstacleField} depicts a 2D obstacle field.



\section{Methods}

\lipsum[1]

\paragraph*{Special paragraph heading}
\lipsum[1]

An article: \cite{WSJrobotRescue}.



\gradingdescdata{0.25}{PDF}{first}{first part of assignment}
\gradingdescdata{0.75}{txt}{second}{second part of assignment}

\makeatletter
\gradingrubric{\@gradeWriteupSections}{\@gradeTableRows}
\makeatother




\stepcounter{section}
\bibliographystyle{plainnat}
\bibliography{RBE_assignment_template}


\stepcounter{section}
\printurls


\vspace{6em}
\begin{center}
{\small \sffamily Last update: \today}
\end{center}

\end{document}


